%\needspace{4\baselineskip}
\color{uniblue}{\subsection{Фиксация положения выключателей (ФПВ) }}
\color{black}

\begin{enumerate}[label=\arabic{section}.\arabic{subsection}.\arabic{enumi}, labelsep=4pt, leftmargin=0pt, itemindent=57pt, itemsep=0pt, parsep=0pt] % если что можно крутить itemsep


\item
Логика ФПВ В1(ЛВ) , ФПВ В2(ОВ) , ФПВ В3, ФПВ В1(ЛВ)СН, ФПВ В2(ОВ)СН аналогичны, 
рассмотрим на примере ФПВ В1(ЛВ). 

\item
Ввод в работу ФПВ В1(ЛВ)  и  осуществляется с помощью уставок \textbf{«ФПВ В1(ЛВ)»} . Оперативный 
ввод/вывод ФПВ производится с помощью БУ (блоков управления функциями) «Вывод ФПВ В1(ЛВ) ». 
Предусматривается возможность формирования положения выключателей с контролем 
разъединителей, которая задается уставками \textbf{«Контроль Р В1(ЛВ)»}. С помощью уставок \textbf{«Схема Р 
В1(ЛВ)»} имеется возможность выбрать схему выключателя, два или три разъединителя. С помощью 
уставки \textbf{«Тип В1(ЛВ)»} имеется возможность выбрать тип выключателя, линейный или обходной. Так 
же от блок-контактов В1(ЛВ)   формируется  внутренний пуск ЗНФ, который пофазно контролирует 
расхождение контактов полюсов выключателей, возникающее при трехфазном включении или 
отключении выключателя. 



\small
\setlength{\extrarowheight}{0.05cm} % высота ячеек в таблице
\begin{longtable}
{|>{\centering\arraybackslash}m{4.5cm}|>{\centering\arraybackslash}m{4.5cm}|>{\centering\arraybackslash}m{3cm}|>{\centering\arraybackslash}m{3.5cm}|}
\caption{Перечень уставочных параметров ФПВ\hfill\vspace{-0.5\baselineskip}}\label{ksn:tbl2} \\ 
\hhline{|-|-|-|-|}
\rowcolor{gray!30}
Наименование уставки & Диапазон & Шаг &  Значение по умолчанию \\ 
\hhline{|-|-|-|-|}
\endfirsthead

\caption*{\hspace{3pt}\emph{Продолжение таблицы \ref{ksn:tbl2}\hfill\vspace{-0.5\baselineskip}}} \\ % сделано по ГОСТ 2.105 п.6.8.7
\hhline{|-|-|-|-|}
\rowcolor{gray!30}
Наименование уставки & Диапазон & Шаг &  Значение по умолчанию \\  
\hhline{|-|-|-|-|}
\endhead
\hhline{|-|-|-|-|}

\endfoot
\endlastfoot

%%%%%%%%%%%%содержимое%%%%%%%%%%%%%%

\end{longtable}
\setlength{\extrarowheight}{0.0cm} % возврат параметра высота ячеек в таблице
\normalsize


\begin{figure}[H]
\centering
\includegraphics[width=1\textwidth,height=1\textheight,keepaspectratio]{\fbpath/ФПВ/img/FPV_1.pdf}
\end{figure}

\begin{figure}[H]
\centering
\includegraphics[width=1\textwidth,height=1\textheight,keepaspectratio]{\fbpath/ФПВ/img/FPV_2.pdf}
\end{figure}


\begin{figure}[H]
\centering
\includegraphics[width=1\textwidth,height=1\textheight,keepaspectratio]{\fbpath/ФПВ/img/FPV_logic.pdf}
\captionof{figure}{Функциональная схема ФПВ В1}\label{pict:FPV}
\end{figure}











\end{enumerate} 